\subsubsection{Các biểu đồ so sánh}

\paragraph{Ma trận nhầm lẫn}

\begin{figure}[htbp]
  \centering
  \includegraphics[width=0.9\textwidth]{classification-confusion-matrix}
  \caption{Ma trận nhầm lẫn của các mô hình trên tập kiểm tra}
  \label{fig:classification-confusion-matrix}
\end{figure}

\subsubsection{Phân tích ma trận nhầm lẫn}

Dựa trên các ma trận nhầm lẫn, có thể thấy rõ sự khác biệt trong hành vi
phân loại của các mô hình, đặc biệt trong cách xử lý các lớp thiểu số và
mức độ thiên lệch về các lớp chiếm đa số.

Đối với Softmax, các giá trị trên đường chéo tập trung chủ yếu ở các lớp
lớn, chẳng hạn lớp 1 và lớp 2 với lần lượt 29213 và 45254 mẫu được dự đoán
đúng. Tuy nhiên, nhiều lớp nhỏ bị bỏ sót nghiêm trọng. Ví dụ, lớp 5 chỉ có
25 mẫu được dự đoán đúng nhưng bị nhầm sang lớp 2 tới 1696 mẫu; lớp 6 chỉ
có 206 mẫu đúng nhưng bị nhầm sang lớp 3 tới 2056 mẫu. Điều này cho thấy
mô hình có xu hướng dồn dự đoán về một vài lớp chính, dẫn đến mất cân bằng
trong phân loại.

Ngược lại, LDA cho thấy phân phối dự đoán đồng đều hơn giữa các lớp.
Các lớp thiểu số được nhận diện tốt hơn, chẳng hạn lớp 5 có 419 mẫu được
dự đoán đúng (so với 25 của Softmax) và lớp 6 có 1828 mẫu đúng (so với 206).
Tuy nhiên, vẫn tồn tại nhầm lẫn đáng kể giữa các lớp, ví dụ lớp 3 bị nhầm
sang lớp 6 tới 2371 mẫu. Dù vậy, mô hình không bị hiện tượng tập trung vào
một vài lớp như Softmax, thể hiện sự cân bằng tốt hơn.

Đối với QDA, ma trận nhầm lẫn cho thấy sự phân loại thiếu ổn định với nhiều
giá trị lớn nằm ngoài đường chéo. Ví dụ, lớp 1 bị nhầm sang lớp 5 tới 11997
mẫu và sang lớp 7 tới 18780 mẫu. Tương tự, lớp 2 bị nhầm sang lớp 5 tới
31626 mẫu. Những giá trị này cho thấy mô hình đưa ra nhiều dự đoán sai lệch
nghiêm trọng. Mặc dù một số lớp có Recall cao, Precision lại rất thấp, dẫn
đến hiệu năng tổng thể kém.

Trong khi đó, Weighted Softmax cải thiện đáng kể khả năng nhận diện các lớp
thiểu số. Ví dụ, lớp 5 có 1438 mẫu được dự đoán đúng (so với 25 của Softmax).
Tuy nhiên, điều này đi kèm với việc gia tăng nhầm lẫn giữa các lớp. Chẳng hạn,
lớp 2 bị nhầm sang lớp 5 tới 10934 mẫu và sang lớp 6 tới 2477 mẫu. Ma trận
nhầm lẫn cho thấy các dự đoán bị phân tán rộng hơn, dẫn đến việc tăng Recall
nhưng giảm Precision và Accuracy.

Một điểm chung rõ ràng giữa các mô hình là đều xảy ra nhầm lẫn mạnh giữa
một số cặp lớp cụ thể, đặc biệt là giữa lớp 1 và lớp 2. Cụ thể, ở Softmax,
có 12181 mẫu từ lớp 1 bị nhầm sang lớp 2 và 9793 mẫu theo chiều ngược lại.
Xu hướng này cũng xuất hiện ở LDA (11540 và 10227 mẫu) và Weighted Softmax
(6906 và 12813 mẫu), cho thấy đây là cặp lớp khó phân biệt đối với mọi mô hình.

Ngoài ra, lớp 2 có xu hướng bị nhầm sang nhiều lớp khác ở tất cả các mô hình,
đặc biệt là sang lớp 3, lớp 5 và lớp 6. Ví dụ, ở Weighted Softmax, có 10934 mẫu
từ lớp 2 bị nhầm sang lớp 5 và 2477 mẫu sang lớp 6; ở Softmax và LDA có tương
ứng 1332 và 1013 mẫu bị nhầm từ lớp 2 sang lớp 3. Điều này cho thấy lớp 2 có
phân bố đặc trưng chồng lấn với nhiều lớp khác.

Bên cạnh đó, các lớp thiểu số như lớp 5 đều có số lượng dự đoán đúng thấp
trong hầu hết các mô hình. Ví dụ, Softmax chỉ dự đoán đúng 25 mẫu cho lớp 5,
trong khi LDA và Weighted Softmax cải thiện hơn (419 và 1438 mẫu) nhưng vẫn
thấp so với các lớp lớn.

Những điểm chung này cho thấy các lỗi phân loại mang tính hệ thống, đặc biệt
liên quan đến sự chồng lấn đặc trưng giữa một số lớp và sự mất cân bằng dữ
liệu, thay vì chỉ xuất phát từ từng mô hình riêng lẻ.

Tổng thể, phân tích ma trận nhầm lẫn cho thấy rõ sự đánh đổi giữa các mô hình:
Softmax có xu hướng thiên lệch về lớp lớn, Weighted Softmax mở rộng khả năng
nhận diện nhưng kém chính xác hơn, LDA đạt được sự cân bằng tương đối, trong
khi QDA thể hiện hiệu năng không ổn định và kém hiệu quả.
